%! ~~~ Packages Setup ~~~ 
\documentclass[]{../../../../tex/classes/styledArticle}
\usepackage{../../../../tex/packages/hebrewSupport}
\usepackage{../../../../tex/packages/mathShortcuts}
\usepackage{../../../../tex/packages/theoremsSupport}
\usepackage{../../../../tex/packages/enfancysections}

%! ~~~ Document ~~~

\author{שחר פרץ}
\title{\textit{לינארית 5} $\sim$ משפט קיילי־המילטון}
\date{21 לאפריל 2025}
\begin{document}
	\maketitle
	\section{\en{Reminders}}
	\defi{מטריצה \textit{ניתנת לישלוש} אם היא  דומה למשולשית. }
	\defi{ט"ל $T \co V \to V$ משולשית אם ישנו בסיס $B \subseteq V$ כך ש־$[T]_B$ משולשית. }
	\begin{Theorem}
		\textbf{\textit{(קיילי המילטון)}} תהי $T \co V \to B$ ט"ל ($V$ נ"ס) או $A \in M_n(\F)$, ו־$f_A(x) = f_T(x)$ הפ"א, אז $f_T(T) = 0, \ f_A(A) = 0$. 
	\end{Theorem}
	\theo{אם $A$ מייצגת של העתקה $T$. ו־$f \in \F[x]$, אז $f(A) = 0 \iff f(T) = 0$. }
	(באופן כללי, עבור כל פולינום). 
	\theo{ט"ל משולשית ו־$A$ ניתנת לשילוש, אמ"מ הפ"א האופייני שלהם מתפצל לגורמים לינארים. }
	
	\section{\en{Proof for Cayley–Hamilton Theorem}}
	\begin{proof}
		נוכיח את המשפט בשלושה שלבים – 
	\begin{itemize}
		\item נניח ש־$T$ ניתנת לשילוש. אזי, קיים בסיס $B =: (v_1 \dots v_n)$ כך ש־$[T]_B$ משולשית (עליונה). זאת מתקיים אמ"מ $\forall i \in [n] \co Tv_i \in \Sp(v_1 \dots v_i)$. נפנה להוכיח את משפט קיילי־המילטון למקרה זה. 
		
		\begin{proof}
			\begin{itemize}
				\item \textit{בסיס:} בעבור $n = 1$, אז קיים $\lg \in \F$ כך ש־$f_T(T) = T - \lg I = 0$ (העתקה לינארית חד ממדית היא כפל בסקלר). בפרט $\forall v \in V \co (T - \lg)v = 0$. 
				\item \textit{צעד:} נניח ש־$B = (v_1 \dots v_n, v_{n + 1})$ שעבורו $[T]_B$ משולשית. קיים תמ"ו $W = \Sp(v_1 \dots v_n)$ כך ש־$\dim W \le \dim V$, ועבורו נכון $\forall w \in W \co Tw \in W$ (נכון עבור וקטורי הבסיס, ונכון לכל $w \in W$ מלינאריות). נגדיר $T|_W \co W \to W$ את הצמצום של $T$ ל־$W$. ידוע ש־$T|_w$ ניתנת לשילוש ולכן מקיימת את תנאי האינדוקציה. לכן, $\forall w \in W\co f_{T|_W}(T)(w) = 0$. אזי $f_{T|_W}(x) = \prod_{i = 1}^{n}(x - \lg_i)$ וסה"כ $f_{T}(x) = (x - \lg_{n + 1})f_{T|_W}(x)$ וקיבלנו $\forall w \in W \co f_T(T)(w) = 0$. 
				
				מספיק להראות ש־$\forall v \in V \co (T - \lg_{n + 1})v \in W$. למה? כי: 
				\[ f_T(T)(v) = \cl{\prod_{i = 1}^{n}(T - \lg_i)}(T - \lg_{n + 1})(v) \]
				מלינאריות, מספיק להראות ש־$(T - \lg_{n + 1})(v_{n + 1})\in W$, שכן זה מתקיים על כל בסיס אחר. אך זה ברור – עבור $[T]_B$ העמודה האחרונה היא: 
				\[ \pms{\ag_1 \\ \vdots \\ \ag_n \\ \lg_{n +1}} \]
				ולכן: 
				\[ T(v_{n + 1}) = \ag_1 v_1 + \cdots + \ag_n v_n + \lg_{n + 1} v_{n + 1} \implies (T - \lg_{n + 1})(v_{n + 1}) = \ag_1v_1 + \cdots + \ag_n v_n \in W \]
			\end{itemize}
		\end{proof}
		\item נוכיח בעבור מטריצה משולשית/ניתנת לשילוש. \begin{proof}
			אם $A$ משולשית, אז $f_A(x) = f_{T_A}(x)$ כאשר $T_A \co \F^n \to \F^n$ המוגדרת ע"י $T_A(v) = Av$, ואז $T_A$ ניתנת לשילוש וסיימנו. 
			
			אם $A$ ניתנת לשילוש, אז היא דומה למשולשית, והן בעלות אותו הפולינום האופייני. אז חזרה לתחילת ההוכחה. 
		\end{proof}
		\item עבור $T$ כללית או $A$ כללית. \begin{proof}
			נניח $A = [T]_B$ עבור בסיס $B$, וידוע $f_T(x) = f_A(x)$. ידוע ש־$A$ ניתנת לשילוש אמ"מ $f_A(x)$ מתפצל. \textit{טענה מהעתיד הלא רחוק: }לכל שדה $\F$ קיים שדה $\F \subseteq K$ סגור אלגברית (כל פולינום מעל שדה סגור אלגברית מפתצל). על כן, ניתן לחשוב על $A \in M_n(\F)$ כמו $A \in M_n(K)$. הפולינום האופייני מעל $K$ הוא אותו הפולינום האופייני מעל $\F$. לכן הוא מתפצל (מעל $K$), ולכן הוא דומה למשולשית, ומהמקרה הראשון $f_A(A) = 0$. זאת כי $f_A(A)$ לא תלוי בשדה עליו אנו עובדים, וסה"כ הוכחנו בעבור מטריצה כללית, ולכן לכל ט"ל. 
		\end{proof}
	\end{itemize}
	\end{proof}
	
	\textbf{הערה על שדות סגורים אלגברית. }(אולי לא נאמר בקורס) העובדה שלכל שדה יש שדה שסגור אלגברית – טענה שתלויה באקסיומת הבחירה. הסגור האלגברי הוא יחיד. 
	
	\section{\en{Rings}}
	מה זה אובייקט אלגברי? דוגמאות: תמורות, חבורות, שדות, מרחבים וקטורים, ועוד. הרעיון – "Data עם אקסיומות". 
	
	\defi{\textit{חוג עם יחידה} הוא קבוצה עם שתי פעולות, כפל וחיבור, ניטרלים לפעולות (0, 1) כך שמתקיימות כל אקסיומות השדה למעט (אולי) קיום איבר הופכי, וקומטטיביות הכפל. }
	
	אנחנו נתעניין ספצפית בחוגים קומטטיבים, כלומר, בהם הכפל כן קומטטיביים. המטריצות הריבועיות מעל אותו הגודל, לדוגמה, הוא חוג שאיננו קומטטיבי. החוג ה"בסיסי ביותר" – חוג השלמים (אין הופכי). ישנם חוגים בלי יחידה (לדוגמה הזוגיים בלי יחידה), לא נדבר עליהם. 
	
	\ndoc
\end{document}
